 \documentclass[11pt]{article}
\usepackage[a4paper, top=2.5cm, bottom=2.5cm, left=2.5cm, right=2.5cm]{geometry}
\usepackage{cite}
\usepackage{hyperref}
\usepackage{amsmath}
\usepackage{setspace}
\usepackage{hyperref}
\usepackage{url}
\usepackage{pgfgantt}
\usepackage{titlesec}
\bibliographystyle{IEEEtran}
\usepackage{listings}
\usepackage{xcolor}
\usepackage{attachfile2}
\usepackage{hyperref}

\usepackage{attachfile2}
\usepackage{makecell}

\lstdefinestyle{mystyle}{
  backgroundcolor=\color{gray!5},
  commentstyle=\color{green!50!black},
  keywordstyle=\color{blue!80!black},
  numberstyle=\tiny\color{gray},
  stringstyle=\color{orange!90!black},
  basicstyle=\ttfamily\small,
  breaklines=true,
  numbers=left,
  numbersep=6pt,
  frame=single,
  tabsize=4,
  showstringspaces=false
}

\lstset{style=mystyle}

\titlespacing*{\section}{0pt}{*0.8}{*0.65}
\setstretch{0.85}
\titleformat{\section}{\large\bfseries}{\thesection.}{0.5em}{}

\begin{document}

\begin{center}
{\Large \textbf{Group Assignment}}\\[8pt]
{\small \textbf{COMP 6321 - Fall 2025}}\\[5pt]
{\small Ayda Nodel Hokmabadi, Saumya Rana, Zar Chi Phyo, Mounika Satya Vani Gundavarapu}\\[3pt]
\end{center}

\vspace*{-2em}
\normalfont\normalsize


\begin{abstract}

\end{abstract}

\section{Introduction}


As the Internet and email have grown in popularity over the years, spam has become a major concern for user experience and cyber security. As spam emails become sophisticated, it has a severe impact to financial concerns for businesses and frustrates individual recipients. \cite{shazad2024spam}. Traditional approaches to spam detection, such as rule-based or keyword-matching, are becoming less effective, since scammers can track and manipulate the traditional spam email detection methods, and resulting in to attack these predefined systems. \cite{khan2024adversarial}. To improve classification accuracy and generalization, natural language processing (NLP) and machine learning techniques are being used. There are several techniques to classify the spam email such as Logistic Regression (LR), Support Vector Machines (SVM) \cite{cervantes2020comprehensive}, K-nearest neighbour (KNN) \cite{abu2019effects}, Decision Trees, Random Forests and Naive Bayes using bag-of-words and TF-IDF features. In contrast to the conventional technique, transformer-based models such as BERT by Google, have demonstrated pinpoint accuracy and higher efficiency in real-world applications of text classification tasks \cite{agusman2025comparative}. Therefore, in this project, we aim to implement a spam classification system that compares classical (Naive Bayes (NB) and Logistic Regression (LR) as baseline modelS) and modern learning models (i.e; learning with pre-trained BERT \cite{devlin2019bert} and DistilBERT \cite{sanh2019distilbert}) models. By comparing these methods, the project seeks to evaluate how transfer learning improves spam detection in terms of accuracy, precision, recall, and F1-score, contributing to more effective and context-aware email filtering.

\section{Model Selection}

Deciding which pre-trained model to be used plays an essential role in a project. Traditional methods cannot capture the semantic relationships of the context. Although LSTM models shows high efficiency, they have limitation on training data size while transformers can learn the complicated sentence structure and deep contextual understanding of text. Moreover, transformers adapt well to new patterns and generalize better for unseen data. While considering for models, the trade-off between bias and variance is also fundamental, whether the chosen model has a high bias i.e, too simple to learn the data or a high variance, i.e, too complex patterns. Since classical models like Logistic Regression and Naive Bayes have high influence on data and features and low variance with poor generalization which can lead to underfit the model. In contrast, BERT/DistilBERT model has 12 encoders with 110 million settings \cite{shazad2024spam}, it can easily learn small datasets and ability to capture deep semantic, therefore, a low bias nature but high variance with proper regularization and hyperparameter tuning can mitigate for overfitting.  Therefore, for our spam email classification project, we proposed to use BERT and DistilBERT, which is a lightweight, distilled version of BERT-based model can give low computational complexity while maintaining high performance. In this project, we use the public Bourique/data email spam corpus available on Hugging Face and our system is classified whether the input sample is spam or not by adding a fully connected classification layer on top of the Transformer encoder.



\section{Hyperparameter Tuning}


Through hyperparameter tuning, we tried to increase the performance and the stability of our spam classification models which were based on transformers. The most important hyperparameters that have considerable impact on the training performance of our transformer models (BERT, DistilBERT) \cite{devlin2019bert} are the learning rate, batch size, number of epochs, weight decay, warm-up ratio, and maxium sequence length.
Learning rate is the most sensitive parameter for fine-tuning the transformer and controls the extent of the model weights update . Batch size has an effect on the gradient stability and the speed of training. Epochs decide the number of passes through the data that the model will make; very few will result in underfitting, whereas, too many will lead to overfitting. Weight decay is a moderation technique, which makes the model less likely to memorize the training data. Warm-up ratio is a method of training that helps to prevent any sudden fluctuations in the learning rate by increasing it slowly at the beginning. Max sequence length determines the amount of information retained from each email in terms of context, and at the same time it's also a factor influencing the computational cost.
We employed a tuning strategy that fused together small grid search and iterative manual refinement; this technique is the one generally used for transformer models due to their massive training costs \cite{sanh2019distilbert}. Learning rates (1e-5 to 5e-5), batch sizes (8, 16, 32), and epochs (2-5) were tested and the rest of the settings were kept the same. Validation F1-score was used to evaluate each configuration, which gave us the opportunity to weigh model complexity against generalization performance. The best results were given by lower learning rates and moderate epochs; longer or deeper training caused overfitting without gaining anything in terms of performance. This tuning had a direct impact on the model selection, allowing us to compare BERT and DistilBERT based on their validation performance and choose the configuration that generalized best under the tuned hyperparameters.

\section{Under-Fitting and Over-Fitting}

Under-Fitting is when a model performs poorly on training data and poorly on test data. The model is too simple to learn the patterns in the data. Overfitting, on the other hand, is when a model performs well on trainig data but poorly on test data. In this case, the model has memorized the training set including noise, and therefore, fails to generalize to new, unseen patterns. 


In our project, we observed these concepts through train vs test evaluation metrics. For baseline models (Naive Bayes and Logistic Regression), the training and test accuracies and F1-score were almost the same. For example Logistic Regression achieved 98.36\% training accuracy and 98.14\% test accuracy, while Naive Bayes achieved 95.69\% training accuracy and 95.53\% test accuracy. This very small gap shows that neither it was overfitting or undefitting. Both models generalized well on our TF-IDF features. This make sense because TF-IDF combined with linear models is a high-bias, low-variance setup. These models learn simple decision boundaries which are based on word frequencies, and the spam dataset contains many clear keyword patterns, so these simple models do not overfit easily. When we moved to transformer-based models (BERT and DistilBERT), we observed a different trend. These models achieved extremely high performance on the training data. For example, during fine-tuning BERT's training accuracy approached 100\%, while its test accuracy stabilized arround 99.22\%-99.28\% <Other metrics?}. Although these results are strong, the gap between almost perfect training accuracy and slightly lower test accuracy indicates a higher risk of overfitting. This happens because transformer models contain millions of parameters and can easily memorize the training set if tuning is not handled properly. Without appropriate choices of learning rate, batch size, number of epochs, these models tend to fit noise rather than learn generalizable patterns.

These observation also connect directly to hyperparameter tuning and model selection. The baseline TF-IDF models required almost no tuning because they are inherently high-bias, low-variance models. Their simplicity and strong regularization properties prevent them from memorizing the training data, making extensive hyperparameter tuning unnecessary. In contrast, BERT and DistilBERT were highly senstive to hyperparameters (selecting an appropriate learning rate, using small numbers of epochs, and choosing the suitable batch sizes), helped reduce overfitting and stabilized test performance.
Choosing the right model depends on balancing complexity and generalization as simple models avoid overfitting but cannot capture deeper context and semantic patterns, while more complex transformer-based models have a higher tendency to overfit due to their size. But once properly tuned, they generalize better than baseliness because they are capable of learning richer contexual meaning.

\section{Model Evaluation}
Model evaluation is one of the most important parts of any machine learning project. It shows us how well our model performs on the data we have, and more importantly, how well it will work on new data it hasn't seen before \cite{hastie2009elements}. Good evaluation means picking the right metrics that actually matter for our specific problem. In classification tasks, just looking at accuracy isn't enough, especially when making different types of mistakes has different consequences \cite{sokolova2009systematic}. That's why we use metrics like precision, recall, and F1-score to get a clearer picture. Precision tells us how many of our positive predictions were right, recall shows us how many of the actual positives we caught, and F1-score combines both to give us a balanced view of performance \cite{powers2020evaluation}. When we look at these metrics alongside training versus test results, we can see whether our model is under-fitting, over-fitting, or finding the right balance \cite{japkowicz2011evaluating}.  In our spam detection project, proper evaluation was essential. For our baseline models—Naive Bayes and Logistic Regression with TF-IDF features, we noticed that accuracy and F1-scores were almost the same on both training and test data. This told us the models were learning well and not just memorizing the training examples. With the transformer models like BERT and DistilBERT, things looked a bit different. Training accuracy reached nearly 100\%, but test performance was slightly lower. This gap, even though it was small, showed us that these bigger models were starting to memorize the training data instead of learning patterns that work everywhere \cite{devlin2019bert, sanh2019distilbert}. By watching the validation F1-score during training and adjusting the learning rate and number of epochs, we managed to reduce this problem and get more reliable results. In the end, careful evaluation helped us compare all the models fairly and pick the ones that gave us the best mix of strong performance and dependability, which is exactly what we need for real spam filtering \cite{shazad2024spam}. 

\section{Learning with Unbalanced Data}

The unbalanced or imbalanced data is defined as one type of data is considerably larger than that of other in the dataset, i.e, the proportion of majority class (ham) is significantly higher than minor class (spam). When the model predicts the data, it priotize on major class and consider the minor as noise and ignore it, resulting in high accuracy of overall system but low recall and predition in spam classification. When there is significant imbalance, accuracy becomes unreliable. Therefore, learning with unbalanced data can lead to several major impacts directly on model performance and even poor predictions on the dataset \cite{werner2023imbalanced}. To mitigate these issues, numerous techniques are used such as undersampling, oversampling, reweighting and synthetic data generation (SMOTE (Synthetic Minority Over-sampling Technique)). 

\section{Discussion and Conclusion}



 
\newpage
\bibliography{references}




\end{document}

